\section{\texttt{@api.model} and \texttt{@api.model\_create\_multi}}

\begin{frame}[fragile]{\texttt{@api.model}}
	\begin{itemize}
		\item Marks a method as a \textbf{model-level} method.
		\item \texttt{self} is an empty recordset of the model, not one specific record.
\begin{block}{Method Syntax}
\begin{lstlisting}
@api.model
def my_method(self, values):
	return self.create(values)
\end{lstlisting}
\end{block}
		\item Common for helpers that create/search without needing an existing record.
	\end{itemize}
\end{frame}

\begin{frame}[fragile]{Breaking Down \texttt{@api.model}}
\begin{lstlisting}
@api.model
def create_demo_student(self):
	print(self)  # student.student()
	return self.create({'name': 'Demo Student'})
\end{lstlisting}
	\begin{itemize}
		\item Call it from the model:
\begin{lstlisting}
self.env['student.student'].create_demo_student()
\end{lstlisting}
		\item Do not expect \texttt{self.id} to exist inside a pure model method.
	\end{itemize}
\end{frame}

\begin{frame}[fragile]{\texttt{@api.model\_create\_multi}}
	\begin{itemize}
		\item Used for efficient multi-record creation.
		\item The argument is a \textbf{list of dictionaries}.
\begin{lstlisting}
@api.model_create_multi
def create(self, vals_list):
	for vals in vals_list:
		vals.setdefault('active', True)
	return super().create(vals_list)
\end{lstlisting}
		\item Preferred modern override for \texttt{create()} in Odoo 17.
	\end{itemize}
\end{frame}

\begin{frame}[fragile]{\texttt{@api.model} vs Record Methods}
	\begin{itemize}
		\item Model method:
\begin{lstlisting}
@api.model
def get_active_count(self):
	return self.search_count([('active', '=', True)])
\end{lstlisting}
		\item Record method:
\begin{lstlisting}
def action_archive(self):
	self.ensure_one()
	self.active = False
\end{lstlisting}
		\item Choose based on whether you need existing records.
	\end{itemize}
\end{frame}
